% answers_short.tex —— 2016 年·北京卷·文科 · 简短答案速查
% 由 tools/make_short_answers.py 自 answers.tex 的 \gksolhead/\gkanswer 提取，
% 并与 2016/2016_数学真题_北京文_答案要点.md 逐题比对后生成，勿手改（改源 MD 后重跑）。

\begin{gknotice}
\textbf{使用说明}：本节只给最终答案便于快速核对，完整推导见解析版。答案为编写者独立推导，\textbf{非官方标准答案}。
\end{gknotice}

\begin{gknotice}
选择题答案中的 U 表示待核，未给出猜测答案，详情见解析版。
\end{gknotice}

\gkpart{选择题}
\begin{center}\begin{tabular}{ccccccccc}
\toprule
题号 & 1 & 2 & 3 & 4 & 5 & 6 & 7 & 8 \\
\midrule
答案 & C & A & B & D & C & B & C & U \\
\bottomrule
\end{tabular}\end{center}

\gkpart{填空题}
\begin{itemize}[leftmargin=2.6em]
  \item[\textbf{9.}] \(\dfrac{\pi}{6}\)
  \item[\textbf{10.}] \(2\)
  \item[\textbf{11.}] \(\dfrac{3}{2}\)
  \item[\textbf{12.}] \(a=1\)；\(b=2\)。
  \item[\textbf{13.}] \(1\)
  \item[\textbf{14.}] ① \(16\)；② \(29\)。
\end{itemize}

\gkpart{解答题}
\begin{itemize}[leftmargin=2.6em]
  \item[\textbf{15.}] (1) \(a_n=2n-1\)；(2) \(S_n=n^{2}+\dfrac{3^{n}-1}{2}\)
  \item[\textbf{16.}] (1) \(\omega=1\)；(2) \(\left[k\pi-\dfrac{3\pi}{8},\ k\pi+\dfrac{\pi}{8}\right]\)（\(k\in\mathbf{Z}\)）
  \item[\textbf{17.}] (1) \(w=3\)；(2) \(10.5\) 元
  \item[\textbf{18.}] (1) 见解析版；(2) 见解析版；(3) 存在，\(F\) 为棱 \(PB\) 的中点
  \item[\textbf{19.}] (1) \(\dfrac{x^{2}}{4}+y^{2}=1\)，离心率 \(e=\dfrac{\sqrt{3}}{2}\)；(2) 见解析版
  \item[\textbf{20.}] (1) \(y=bx+c\)；(2) \(0<c<\dfrac{32}{27}\)；(3) 见解析版
\end{itemize}
